% Copyright (c) 2026 Selina Cheggour. Tous droits réservés. Voir le fichier LICENSE.
% Lexique simple et pistes pour aller plus loin.

\bandeaupartie{Lexique et pistes}{Les mots du TP, en une phrase chacun ; puis ce qu'on pourrait faire ensuite.}

\section{Lexique}

\renewcommand{\arraystretch}{1.32}
\rowcolors{2}{notebg}{white}
\begin{tabularx}{\linewidth}{>{\bfseries\leavevmode\color{darkblue}}p{3.7cm}X}
\rowcolor{titlebg}\color{white}mot & \leavevmode\color{white}\textbf{en une phrase} \\
verbatim & le texte libre écrit par le consommateur, tel qu'il est arrivé. \\
apprentissage non supervisé & on cherche une structure dans des données sans étiquette : ici, des groupes de plaintes. \\
TF-IDF & une coordonnée par mot ; un mot pèse lourd s'il est fréquent dans la plainte et rare ailleurs. \\
cosinus & la proximité de deux vecteurs : 1 s'ils pointent dans la même direction, 0 s'ils n'ont rien en commun. \\
jeton & le morceau de texte que lit un transformer : un mot fréquent, ou un bout de mot rare. \\
transformer & un réseau de neurones où chaque jeton \og regarde \fg{} tous les autres (l'attention) pour préciser son sens. \\
plongement & le vecteur de sens qu'un transformer donne à un texte entier. \\
centrage par langue & retirer à chaque vecteur la moyenne de sa langue, pour ne garder que ce que la plainte dit. \\
t-SNE & une méthode qui dessine en 2 ou 3 dimensions des points proches dans un grand espace ; les axes n'ont pas de sens. \\
ACP & l'analyse en composantes principales : garder les quelques directions où les données varient le plus. \\
HDBSCAN & un regroupement par densité, qui trouve seul le nombre de groupes et laisse de côté les points isolés. \\
Ward & une classification hiérarchique qui range chaque point, en fusionnant les groupes qui augmentent le moins la dispersion. \\
groupes jumeaux & le même problème trouvé deux fois, une fois par langue. \\
\end{tabularx}

\begin{tabularx}{\linewidth}{>{\bfseries\leavevmode\color{darkblue}}p{3.7cm}X}
\rowcolor{titlebg}\color{white}mot & \leavevmode\color{white}\textbf{en une phrase (suite)} \\
c-TF-IDF & le TF-IDF calculé sur un groupe entier, pour trouver ses mots les plus caractéristiques. \\
LLM & un grand modèle de langue qui écrit du texte ; ici, il tourne dans le notebook et propose titres et définitions. \\
ARI & l'indice de Rand ajusté : l'accord entre deux regroupements, 1 s'ils sont identiques, 0 au hasard. \\
loi de Poisson & la loi du nombre d'événements rares quand on en attend en moyenne $\lambda$. \\
p-valeur & la probabilité d'observer au moins autant si le hasard seul jouait. \\
Benjamini-Hochberg & une correction des p-valeurs quand on fait beaucoup de tests à la fois. \\
part dans le pic & la part des plaintes d'un problème, sur une période, qui vient d'une seule marque. \\
quasi doublon & une plainte dont la plus proche voisine TF-IDF a un cosinus supérieur à 0,9. \\
incident candidat & un problème en excès sur une marque et une période, porté surtout par cette marque. \\
pic partagé & un problème en excès sur une période chez plusieurs marques à la fois. \\
campagne & un groupe de messages presque identiques : une mobilisation, pas un défaut. \\
rappel, précision & la part d'un groupe caché retrouvée ensemble ; la part d'un groupe trouvé qui vient du groupe caché. \\
\end{tabularx}

\clearpage
\section{Pour aller plus loin}

\begin{tcbraster}[raster columns=2, raster equal height, raster column skip=8pt, raster row skip=8pt]
\begin{conceptbox}[title={\faLanguage\ Traduire avant de comparer}, unbreakable]
\small Au lieu de centrer par langue, traduire toutes les plaintes dans une seule langue avec un modèle de
traduction qui tourne dans le notebook. Plus lourd, mais sans le risque de retirer une partie du sens avec
la langue.
\end{conceptbox}
\begin{conceptbox}[title={\faCalendar*\ La date de l'événement}, unbreakable]
\small Une plainte peut arriver des semaines après l'achat. Si l'entreprise fournit la date de fabrication
ou le numéro de lot, l'alerte pourrait porter sur le lot plutôt que sur la période de la plainte.
\end{conceptbox}
\begin{conceptbox}[title={\faStore\ D'autres cases que la marque}, unbreakable]
\small Quand une marque domine la base, comparer les marques n'a plus de sens : on peut tester l'excès par
segment, par pays (si la colonne existe), ou par période seule.
\end{conceptbox}
\begin{conceptbox}[title={\faUserCheck\ Une évaluation humaine}, unbreakable]
\small Faire relire, par une personne habilitée, quelques plaintes de chaque problème et noter si le titre
proposé par le LLM est juste : c'est la seule façon de mesurer la qualité sur la base réelle.
\end{conceptbox}
\begin{conceptbox}[title={\faProjectDiagram\ D'autres regroupements}, unbreakable]
\small Réduire les vecteurs avec UMAP avant HDBSCAN, ou essayer une modélisation de sujets complète
(BERTopic), et les ajouter à la chaîne de comparaison, à côté des huit outils.
\end{conceptbox}
\begin{conceptbox}[title={\faBell\ Une alerte au fil de l'eau}, unbreakable]
\small Recalculer chaque semaine les cases anormales sur les nouvelles plaintes seulement, et n'alerter
que sur ce qui change : c'est ainsi qu'un tel outil servirait au quotidien.
\end{conceptbox}
\end{tcbraster}

\vfill
\begin{center}\color{slate}\small
\faUserTie\quad Une question ? Lève la main, l'encadrante passe.
\end{center}
